% Data Mining - Chapter 4: Data Classification
\documentclass[10pt]{article}
% ============================================================
% _preamble.tex -- Shared preamble for all DM chapter docs
% Usage: % ============================================================
% _preamble.tex -- Shared preamble for all DM chapter docs
% Usage: % ============================================================
% _preamble.tex -- Shared preamble for all DM chapter docs
% Usage: \input{../_preamble} AFTER \documentclass, BEFORE \begin{document}
% ============================================================

\usepackage{pslatex}                    % Times New Roman
\usepackage[
    paperheight=22.86cm,
    paperwidth=15.24cm,
    left=1.27cm, right=1.27cm,
    top=1.6cm,   bottom=1.6cm,
    headheight=1in
]{geometry}
\usepackage[utf8]{inputenc}
\usepackage[vietnamese]{babel}
\usepackage[hidelinks]{hyperref}
\usepackage{graphicx}
\usepackage{enumerate}
\usepackage{float}
\usepackage{longtable}
\usepackage{array}
\usepackage{setspace}
\usepackage{sectsty}
\usepackage[normalem]{ulem}
\usepackage{caption}
\usepackage{titlesec}
\usepackage{amsmath}
\usepackage{amssymb}
\usepackage{enumitem}

% ----- Typography settings -----
\sectionfont{\fontsize{12pt}{12pt}\selectfont}
\captionsetup{font={normalsize,rm}, labelfont=bf, skip=10pt}
\titlelabel{\thetitle.\quad}
\urlstyle{same}

% ----- Paragraph indent -----
\usepackage{indentfirst}        % indent the first paragraph of every section
\setlength{\parindent}{1.5em}   % indent width (~4 spaces at 10pt)
\setlength{\parskip}{0pt}       % no extra space between paragraphs
 AFTER \documentclass, BEFORE \begin{document}
% ============================================================

\usepackage{pslatex}                    % Times New Roman
\usepackage[
    paperheight=22.86cm,
    paperwidth=15.24cm,
    left=1.27cm, right=1.27cm,
    top=1.6cm,   bottom=1.6cm,
    headheight=1in
]{geometry}
\usepackage[utf8]{inputenc}
\usepackage[vietnamese]{babel}
\usepackage[hidelinks]{hyperref}
\usepackage{graphicx}
\usepackage{enumerate}
\usepackage{float}
\usepackage{longtable}
\usepackage{array}
\usepackage{setspace}
\usepackage{sectsty}
\usepackage[normalem]{ulem}
\usepackage{caption}
\usepackage{titlesec}
\usepackage{amsmath}
\usepackage{amssymb}
\usepackage{enumitem}

% ----- Typography settings -----
\sectionfont{\fontsize{12pt}{12pt}\selectfont}
\captionsetup{font={normalsize,rm}, labelfont=bf, skip=10pt}
\titlelabel{\thetitle.\quad}
\urlstyle{same}

% ----- Paragraph indent -----
\usepackage{indentfirst}        % indent the first paragraph of every section
\setlength{\parindent}{1.5em}   % indent width (~4 spaces at 10pt)
\setlength{\parskip}{0pt}       % no extra space between paragraphs
 AFTER \documentclass, BEFORE \begin{document}
% ============================================================

\usepackage{pslatex}                    % Times New Roman
\usepackage[
    paperheight=22.86cm,
    paperwidth=15.24cm,
    left=1.27cm, right=1.27cm,
    top=1.6cm,   bottom=1.6cm,
    headheight=1in
]{geometry}
\usepackage[utf8]{inputenc}
\usepackage[vietnamese]{babel}
\usepackage[hidelinks]{hyperref}
\usepackage{graphicx}
\usepackage{enumerate}
\usepackage{float}
\usepackage{longtable}
\usepackage{array}
\usepackage{setspace}
\usepackage{sectsty}
\usepackage[normalem]{ulem}
\usepackage{caption}
\usepackage{titlesec}
\usepackage{amsmath}
\usepackage{amssymb}
\usepackage{enumitem}

% ----- Typography settings -----
\sectionfont{\fontsize{12pt}{12pt}\selectfont}
\captionsetup{font={normalsize,rm}, labelfont=bf, skip=10pt}
\titlelabel{\thetitle.\quad}
\urlstyle{same}

% ----- Paragraph indent -----
\usepackage{indentfirst}        % indent the first paragraph of every section
\setlength{\parindent}{1.5em}   % indent width (~4 spaces at 10pt)
\setlength{\parskip}{0pt}       % no extra space between paragraphs


\begin{document}
\onehalfspacing
\renewcommand{\arraystretch}{1.5}

\newcommand{\ChapterNum}{4}
\newcommand{\ChapterTitle}{PHÂN LỚP DỮ LIỆU (DATA CLASSIFICATION)}
% ============================================================
% _title.tex -- Shared title block for all DM chapter docs
% Required macros (define BEFORE % ============================================================
% _title.tex -- Shared title block for all DM chapter docs
% Required macros (define BEFORE % ============================================================
% _title.tex -- Shared title block for all DM chapter docs
% Required macros (define BEFORE \input{../_title}):
%   \ChapterNum    e.g. {5}
%   \ChapterTitle  e.g. {PHÂN CỤM DỮ LIỆU (DATA CLUSTERING)}
% ============================================================

\begin{center}
    {\fontsize{14pt}{14pt}\selectfont
        \textbf{KHAI PHÁ DỮ LIỆU -- CHƯƠNG \ChapterNum}\par}
    \vspace{4pt}
    {\fontsize{13pt}{13pt}\selectfont
        \textbf{\ChapterTitle}\par}
    \vspace{6pt}
    {\fontsize{10pt}{10pt}\selectfont
        Tổng hợp kiến thức, câu hỏi tự luận và trắc nghiệm\par}
\end{center}
):
%   \ChapterNum    e.g. {5}
%   \ChapterTitle  e.g. {PHÂN CỤM DỮ LIỆU (DATA CLUSTERING)}
% ============================================================

\begin{center}
    {\fontsize{14pt}{14pt}\selectfont
        \textbf{KHAI PHÁ DỮ LIỆU -- CHƯƠNG \ChapterNum}\par}
    \vspace{4pt}
    {\fontsize{13pt}{13pt}\selectfont
        \textbf{\ChapterTitle}\par}
    \vspace{6pt}
    {\fontsize{10pt}{10pt}\selectfont
        Tổng hợp kiến thức, câu hỏi tự luận và trắc nghiệm\par}
\end{center}
):
%   \ChapterNum    e.g. {5}
%   \ChapterTitle  e.g. {PHÂN CỤM DỮ LIỆU (DATA CLUSTERING)}
% ============================================================

\begin{center}
    {\fontsize{14pt}{14pt}\selectfont
        \textbf{KHAI PHÁ DỮ LIỆU -- CHƯƠNG \ChapterNum}\par}
    \vspace{4pt}
    {\fontsize{13pt}{13pt}\selectfont
        \textbf{\ChapterTitle}\par}
    \vspace{6pt}
    {\fontsize{10pt}{10pt}\selectfont
        Tổng hợp kiến thức, câu hỏi tự luận và trắc nghiệm\par}
\end{center}


% ============================================================
\section{TỔNG QUAN VỀ PHÂN LỚP}
% ============================================================

\subsection{Định nghĩa và quy trình}

\textbf{Phân lớp (Classification)} là phương pháp phân tích dữ liệu nhằm trích xuất mô hình mô tả các lớp dữ liệu. Đây là quy trình \textbf{hai bước}:
\begin{enumerate}
    \item \textbf{Huấn luyện (Training):} Xây dựng bộ phân lớp bằng cách học (analyzing) tập huấn luyện có nhãn $\langle X, y \rangle$.
    \item \textbf{Phân lớp (Classification):} Phân lớp các mẫu/đối tượng mới. Nếu độ chính xác của bộ phân lớp chấp nhận được, dùng nó để phân lớp dữ liệu mới.
\end{enumerate}

\textbf{Đặc điểm:} Phân lớp là học \textbf{có giám sát (supervised learning)} -- có nhãn lớp trong dữ liệu huấn luyện.

\textbf{Ví dụ tình huống:}
\begin{itemize}
    \item Email: ``spam'' hay ``not spam''.
    \item Giao dịch trực tuyến: ``gian lận'' hay ``bình thường''.
    \item Y tế: u bướu ``lành tính'' hay ``ác tính''; bệnh nhân ``mắc bệnh'' hay ``không mắc bệnh''.
\end{itemize}

\textbf{Các thuật toán phân lớp phổ biến:}
\begin{itemize}
    \item Logistic Regression, Decision Tree, Bayesian method, ANN, K-NN, Case-based reasoning, Genetic algorithms, Rough/Fuzzy sets.
\end{itemize}

% ============================================================
\section{HỒI QUY LOGISTIC (LOGISTIC REGRESSION)}
% ============================================================

\subsection{Hàm Sigmoid}

Hồi quy tuyến tính $h_\theta(X) = \theta^TX$ có thể cho kết quả $> 1$ hoặc $< 0$, không phù hợp cho phân lớp. Hồi quy Logistic đảm bảo $0 \leq h_\theta(X) \leq 1$ bằng hàm Sigmoid (Logistic function):

\[
h_\theta(X) = g(\theta^TX) = \frac{1}{1 + e^{-\theta^TX}}
\]
\[
g(z) = \frac{1}{1 + e^{-z}}, \quad g(z) \in (0, 1)
\]

\textbf{Ý nghĩa:} $h_\theta(x) = P(y=1 \mid x; \theta)$ -- xác suất nhãn là 1 khi đầu vào là $x$.

\subsection{Decision Boundary (Ranh giới quyết định)}

\begin{itemize}
    \item Dự đoán $y = 1$ khi $h_\theta(X) \geq 0.5 \Leftrightarrow \theta^TX \geq 0$.
    \item Dự đoán $y = 0$ khi $h_\theta(X) < 0.5 \Leftrightarrow \theta^TX < 0$.
\end{itemize}

Decision boundary là tập $\{\theta^TX = 0\}$ -- có thể là đường thẳng (tuyến tính) hoặc đường cong (phi tuyến) tùy thuộc vào đặc trưng.

\subsection{Hàm chi phí và Gradient Descent}

Hàm chi phí logistic:
\[
\text{cost}(h_\theta(x), y) = \begin{cases}
-\log(h_\theta(x)) & \text{nếu } y = 1 \\
-\log(1 - h_\theta(x)) & \text{nếu } y = 0
\end{cases}
\]

Rút gọn: $\text{cost}(h_\theta(x), y) = -y\log(h_\theta(x)) - (1-y)\log(1-h_\theta(x))$

\[
J(\theta) = \frac{1}{N}\sum_{i=1}^{N}\left[-y^{(i)}\log(h_\theta(x^{(i)})) - (1-y^{(i)})\log(1-h_\theta(x^{(i)}))\right]
\]

Tối thiểu hóa $J(\theta)$ bằng gradient descent (dạng công thức giống hồi quy tuyến tính sau khi lấy đạo hàm).

\subsection{Phân lớp đa lớp -- One-vs-Rest}

Với $k$ lớp, huấn luyện $k$ bộ phân lớp $h_\theta^{(i)}(x)$ ($i=1,\ldots,k$), mỗi bộ phân biệt lớp $i$ với tất cả lớp còn lại. Dự đoán lớp $i$ có $h_\theta^{(i)}(x)$ lớn nhất.

% ============================================================
\section{CÂY QUYẾT ĐỊNH (DECISION TREE)}
% ============================================================

\subsection{Cấu trúc và ý nghĩa}

\begin{itemize}
    \item \textbf{Internal node (Nút trong):} Kiểm tra một thuộc tính cụ thể.
    \item \textbf{Branch (Nhánh):} Kết quả của phép kiểm tra.
    \item \textbf{Leaf node (Nút lá):} Nhãn lớp.
\end{itemize}

\subsection{Thuật toán xây dựng cây quyết định}

\textbf{Đặc điểm:} Greedy, divide-and-conquer, đệ quy, top-down. Độ phức tạp: $O(n \times |D| \times \log|D|)$.

\textbf{Các thuật toán:} ID3, C4.5, CART (Classification and Regression Trees -- cây nhị phân).

\subsection{Tiêu chí chọn thuộc tính phân tách}

\paragraph{(a) Information Gain (ID3):}
\[
\text{Info}(D) = -\sum_{i=1}^{m} p_i \log_2(p_i), \quad p_i = |C_{i,D}|/|D|
\]
\[
\text{Info}_A(D) = \sum_{j=1}^{v} \frac{|D_j|}{|D|} \text{Info}(D_j)
\]
\[
\text{Gain}(A) = \text{Info}(D) - \text{Info}_A(D)
\]
Chọn thuộc tính có \textbf{Gain lớn nhất}. Hạn chế: ưu tiên thuộc tính có nhiều giá trị (tạo nhiều partition nhỏ).

\paragraph{(b) Gain Ratio (C4.5):}
\[
\text{SplitInfo}_A(D) = -\sum_{j=1}^{v} \frac{|D_j|}{|D|} \log_2\frac{|D_j|}{|D|}
\]
\[
\text{GainRatio}(A) = \frac{\text{Gain}(A)}{\text{SplitInfo}_A(D)}
\]
Chuẩn hóa Information Gain theo độ phân tách; giảm thiên vị với thuộc tính nhiều giá trị.

\paragraph{(c) Gini Index (CART):}
\[
\text{Gini}(D) = 1 - \sum_{i=1}^{m} p_i^2
\]
\[
\text{Gini}_A(D) = \frac{|D_1|}{|D|}\text{Gini}(D_1) + \frac{|D_2|}{|D|}\text{Gini}(D_2)
\]
Chia nhị phân; chọn thuộc tính và ngưỡng phân tách có \textbf{Gini nhỏ nhất}.

% ============================================================
\section{PHÂN LỚP BAYESIAN}
% ============================================================

\subsection{Định lý Bayes}

\[
P(H \mid X) = \frac{P(X \mid H) \cdot P(H)}{P(X)}
\]

\begin{itemize}
    \item $P(H \mid X)$: xác suất hậu nghiệm (posterior) -- xác suất $X$ thuộc lớp $H$.
    \item $P(X \mid H)$: likelihood -- xác suất quan sát $X$ nếu lớp là $H$.
    \item $P(H)$: xác suất tiên nghiệm của lớp (prior).
    \item $P(X)$: hằng số chuẩn hóa.
\end{itemize}

\subsection{Naive Bayesian Classification}

Giả định \textbf{class conditional independence}: các thuộc tính độc lập với nhau trong cùng một lớp:
\[
P(X \mid C_i) = \prod_{k=1}^{n} P(x_k \mid C_i)
\]

\textbf{Phân lớp:} Chọn lớp $C_i$ có $P(C_i \mid X) \propto P(X \mid C_i) \cdot P(C_i)$ lớn nhất.

\textbf{Tính $P(x_k \mid C_i)$:}
\begin{itemize}
    \item Thuộc tính danh mục: $P(x_k \mid C_i) = \frac{|\{X' \mid x'_k = x_k \wedge X' \in C_i\}|}{|C_{i,D}|}$
    \item Thuộc tính liên tục (Gauss): $P(x_k \mid C_i) = \frac{1}{\sqrt{2\pi}\sigma_i} \exp\left(-\frac{(x_k - \mu_i)^2}{2\sigma_i^2}\right)$
\end{itemize}

\textbf{Vấn đề xác suất bằng 0:} Nếu $P(x_k \mid C_i) = 0$ thì $P(X \mid C_i) = 0$, gây ra ``zero-frequency problem''.

\textbf{Giải pháp Laplace smoothing:}
\[
P(x_k \mid C_i) = \frac{|\{X' \mid x'_k = x_k \wedge X' \in C_i\}| + 1}{|C_{i,D}| + m}
\]
trong đó $m$ là số giá trị khác nhau của thuộc tính $A_k$.

\textbf{Ưu/nhược điểm của Naive Bayes:}
\begin{itemize}
    \item Ưu: Dễ cài đặt, học nhanh, dễ hiểu kết quả, hiệu quả trong nhiều trường hợp.
    \item Nhược: Giả định độc lập lớp điều kiện có thể không thỏa trong thực tế.
\end{itemize}

% ============================================================
\section{MẠNG NƠRON NHÂN TẠO (ANN)}
% ============================================================

\subsection{Mô hình Neuron -- Logistic Unit}

Mỗi neuron nhân tạo là một logistic unit với hàm kích hoạt Sigmoid $g(z)$:
\[
a^{(l)}_j = g\left(\sum_i \Theta^{(l-1)}_{ji} a^{(l-1)}_i\right)
\]

\subsection{Kiến trúc ANN}

\begin{itemize}
    \item \textbf{Layer 1:} Input layer (lớp đầu vào).
    \item \textbf{Layer 2, ..., L-1:} Hidden layers (lớp ẩn).
    \item \textbf{Layer L:} Output layer (lớp đầu ra).
\end{itemize}

Ma trận trọng số $\Theta^{(j)}$ có kích thước $s_{j+1} \times (s_j + 1)$ (với $s_j$ là số nút tại lớp $j$).

\subsection{Feedforward (Forward Propagation)}

Tính giá trị kích hoạt từ lớp đầu vào đến lớp đầu ra:
\[
z^{(l+1)} = \Theta^{(l)} a^{(l)}, \quad a^{(l+1)} = g(z^{(l+1)})
\]
(thêm bias node $a_0^{(l)} = 1$ tại mỗi lớp)

\subsection{Hàm chi phí ANN}

Phân lớp nhị phân ($K=1$):
\[
J(\Theta) = -\frac{1}{N}\sum_{i=1}^{N}\left[y^{(i)}\log h_\Theta(x^{(i)}) + (1-y^{(i)})\log(1-h_\Theta(x^{(i)}))\right]
\]

\subsection{Backpropagation}

Thuật toán tính gradient để tối ưu hóa $J(\Theta)$:
\begin{enumerate}
    \item \textbf{Feedforward:} Tính $a^{(l)}$ cho tất cả các lớp.
    \item \textbf{Tính lỗi:} $\delta^{(L)} = a^{(L)} - y$.
    \item \textbf{Lan truyền ngược:} $\delta^{(l)} = (\Theta^{(l)})^T\delta^{(l+1)} \cdot g'(z^{(l)})$.
    \item \textbf{Tính gradient:} $\frac{\partial J}{\partial \Theta^{(l)}_{ij}} = a^{(l)}_j \cdot \delta^{(l+1)}_i$.
    \item Cập nhật $\Theta^{(l)}$ bằng gradient descent.
\end{enumerate}

% ============================================================
\section{K-NEAREST NEIGHBOR (K-NN)}
% ============================================================

\textbf{Ý tưởng:} Phân lớp đối tượng $X$ bằng cách bỏ phiếu đa số (majority vote) từ $k$ đối tượng gần nhất trong tập huấn luyện.

\textbf{Độ đo khoảng cách:} Thường dùng Euclidean: $d(p, q) = \sqrt{\sum_i (p_i - q_i)^2}$

\textbf{Chọn $k$:}
\begin{itemize}
    \item $k$ quá nhỏ: nhạy cảm với nhiễu.
    \item $k$ quá lớn: có thể chọn đối tượng từ lớp khác.
    \item Gợi ý: $k \leq \sqrt{|D|}$.
\end{itemize}

% ============================================================
\section{ĐÁNH GIÁ VÀ CHỌN MÔ HÌNH PHÂN LỚP}
% ============================================================

\subsection{Tiêu chí đánh giá}

\begin{itemize}
    \item \textbf{Accuracy:} Khả năng nhận dạng đúng các đối tượng.
    \item \textbf{Speed:} Chi phí tính toán khi huấn luyện và sử dụng.
    \item \textbf{Robustness:} Hoạt động tốt với dữ liệu nhiễu hoặc thiếu.
    \item \textbf{Scalability:} Xây dựng và cập nhật bộ phân lớp với tập dữ liệu rất lớn.
    \item \textbf{Interpretability:} Khả năng hiểu cách bộ phân lớp hoạt động.
\end{itemize}

\subsection{Precision, Recall và F-score}

\begin{center}
\begin{tabular}{|c|c|c|}
\hline
 & \textbf{Thực tế: X} & \textbf{Thực tế: !X} \\
\hline
\textbf{Dự đoán: X} & TP & FP \\
\hline
\textbf{Dự đoán: !X} & FN & TN \\
\hline
\end{tabular}
\end{center}

\[
\text{Precision} = \frac{TP}{TP + FP}, \quad \text{Recall} = \frac{TP}{TP + FN}
\]
\[
\text{F-score} = \frac{2 \times P \times R}{P + R}
\]

\subsection{Phương pháp đánh giá}

\begin{itemize}
    \item \textbf{Holdout method:} Chia ngẫu nhiên $D$ thành tập train (2/3) và test (1/3).
    \item \textbf{K-fold cross-validation:} Chia $D$ thành $k$ phần bằng nhau; lặp $k$ lần (dùng phần $i$ để test, phần còn lại để train); tính trung bình.
\end{itemize}

% ============================================================
\section{TÓM TẮT}
% ============================================================

\begin{itemize}
    \item \textbf{Phân lớp:} Học có giám sát; 2 bước: training + classification.
    \item \textbf{Logistic Regression:} Sigmoid function; decision boundary; gradient descent; one-vs-rest đa lớp.
    \item \textbf{Decision Tree (ID3/C4.5/CART):} Tiêu chí phân tách: Information Gain, Gain Ratio, Gini Index.
    \item \textbf{Naive Bayes:} Định lý Bayes + giả định độc lập; Laplace smoothing cho zero-frequency.
    \item \textbf{ANN:} Feedforward + Backpropagation; xử lý phi tuyến phức tạp.
    \item \textbf{K-NN:} Lazy learning, dựa trên khoảng cách.
    \item \textbf{Đánh giá:} Accuracy, Precision, Recall, F-score; Holdout, K-fold cross-validation.
\end{itemize}

% ============================================================
\newpage
\section{CÂU HỎI TỰ LUẬN}
% ============================================================

\begin{enumerate}[leftmargin=*, label=\textbf{Câu \arabic*.}]

    \item Phân lớp dữ liệu (Classification) là gì? Mô tả quy trình 2 bước của phân lớp và giải thích tại sao phân lớp là học có giám sát. So sánh với phân cụm (Clustering).

    \item Giải thích tại sao hồi quy tuyến tính $h_\theta(X) = \theta^TX$ không phù hợp cho bài toán phân lớp. Hàm Sigmoid giải quyết vấn đề này như thế nào?

    \item Trình bày hàm Sigmoid $g(z)$ và ý nghĩa của $h_\theta(x) = P(y=1 \mid x; \theta)$ trong hồi quy Logistic. Quy tắc quyết định phân lớp dựa trên $h_\theta(x)$ và $\theta^TX$ như thế nào?

    \item Giải thích khái niệm \textbf{Decision Boundary} trong hồi quy Logistic. Khi nào thì ranh giới quyết định là đường thẳng? Khi nào là đường cong? Cho ví dụ cụ thể.

    \item Trình bày hàm chi phí của hồi quy Logistic. Tại sao không dùng MSE như hồi quy tuyến tính? Giải thích ý nghĩa của từng trường hợp $y=1$ và $y=0$.

    \item Giải thích phương pháp \textbf{One-vs-Rest} để phân lớp đa lớp bằng Logistic Regression. Mô tả quy trình huấn luyện và dự đoán. Hạn chế của phương pháp này?

    \item Mô tả cấu trúc của \textbf{Decision Tree}. Giải thích các khái niệm: internal node, branch, leaf node. Thuật toán xây dựng cây quyết định có đặc điểm gì (greedy, divide-and-conquer)?

    \item Trình bày và so sánh 3 tiêu chí chọn thuộc tính phân tách trong cây quyết định: \textbf{Information Gain}, \textbf{Gain Ratio} và \textbf{Gini Index}. Khi nào Gain Ratio tốt hơn Information Gain?

    \item Tính Information Gain của thuộc tính ``age'' trong tập dữ liệu AllElectronics: $\text{Info}(D) = 0.940$, $\text{Info}_{\text{age}}(D) = 0.694$. Kết quả Gain(age) = 0.246. Giải thích ý nghĩa của con số này.

    \item Trình bày \textbf{Định lý Bayes}. Giải thích các khái niệm: posterior probability, prior probability, likelihood, evidence. Tại sao phân lớp Bayesian dựa trên ``tối đa hóa posterior''?

    \item Giải thích giả định \textbf{class conditional independence} trong Naive Bayes. Khi nào giả định này có thể không thỏa? Hệ quả khi giả định bị vi phạm là gì?

    \item Trình bày \textbf{Laplace smoothing} trong Naive Bayes. Tại sao cần Laplace smoothing? Công thức điều chỉnh như thế nào? Có phương pháp thay thế nào không?

    \item Mô tả kiến trúc \textbf{Mạng Nơron Nhân Tạo (ANN)}. Giải thích: input layer, hidden layers, output layer, weights matrix $\Theta^{(j)}$, kích thước của $\Theta^{(j)}$.

    \item Trình bày thuật toán \textbf{Feedforward (Forward Propagation)} trong ANN. Vì sao cần thêm bias node $a_0^{(l)} = 1$ tại mỗi lớp?

    \item Giải thích thuật toán \textbf{Backpropagation}. Mục đích của backpropagation là gì? Mô tả cách tính $\delta^{(L)}$, $\delta^{(L-1)}$ và gradient $\frac{\partial J}{\partial \Theta^{(l)}_{ij}}$.

    \item Tại sao \textbf{ANN} được sử dụng cho bài toán phân lớp phi tuyến phức tạp mà Logistic Regression không xử lý được? Minh họa bằng ví dụ bài toán XOR.

    \item Giải thích thuật toán \textbf{K-Nearest Neighbor (K-NN)}. Trình bày cách chọn $k$, độ đo khoảng cách và quy tắc bỏ phiếu đa số. Ưu và nhược điểm của K-NN?

    \item Phân biệt \textbf{Precision} và \textbf{Recall}. Trong tình huống phát hiện gian lận thẻ tín dụng, tiêu chí nào quan trọng hơn? Tại sao cần F-score?

    \item So sánh \textbf{Holdout method} và \textbf{K-fold cross-validation} để đánh giá bộ phân lớp. Phương pháp nào đáng tin cậy hơn và tại sao?

    \item So sánh toàn diện 5 phương pháp phân lớp: Logistic Regression, Decision Tree, Naive Bayes, ANN và K-NN theo các tiêu chí: accuracy, speed, robustness, scalability và interpretability.

\end{enumerate}

% ============================================================
\newpage
\section{CÂU HỎI TRẮC NGHIỆM}
% ============================================================

\begin{enumerate}[leftmargin=*, label=\textbf{Câu \arabic*.}]

    \item Phân lớp (Classification) là loại học nào?
    \begin{enumerate}[label=\Alph*.]
        \item Học không giám sát.
        \item Học có giám sát.
        \item Học tăng cường.
        \item Học bán giám sát.
    \end{enumerate}

    \item Hàm Sigmoid $g(z) = \frac{1}{1+e^{-z}}$ có miền giá trị:
    \begin{enumerate}[label=\Alph*.]
        \item $(-\infty, +\infty)$
        \item $[0, 1]$ (nhưng không đạt đúng 0 hoặc 1, ngoại trừ giới hạn).
        \item $[-1, 1]$
        \item $[0.5, 1]$
    \end{enumerate}

    \item Trong Logistic Regression, dự đoán $y = 1$ khi:
    \begin{enumerate}[label=\Alph*.]
        \item $h_\theta(x) \geq 0.5 \Leftrightarrow \theta^Tx \geq 0$.
        \item $h_\theta(x) < 0.5$.
        \item $\theta^Tx < 0$.
        \item $h_\theta(x) = 1$.
    \end{enumerate}

    \item $h_\theta(x) = P(y=1 \mid x; \theta) = 0.7$ có nghĩa là:
    \begin{enumerate}[label=\Alph*.]
        \item $70\%$ đối tượng thuộc lớp 0.
        \item Xác suất $70\%$ đối tượng thuộc lớp 1 với đầu vào $x$.
        \item Mô hình chính xác $70\%$.
        \item $\theta^Tx = 0.7$.
    \end{enumerate}

    \item Hàm chi phí Logistic Regression khi $y=1$ và $h_\theta(x) \to 0$ sẽ:
    \begin{enumerate}[label=\Alph*.]
        \item Tiến về 0 (cost nhỏ).
        \item Tiến về vô cực (cost rất lớn -- phạt mạnh dự đoán sai).
        \item Bằng 0.5.
        \item Bằng $-\log(0.5)$.
    \end{enumerate}

    \item Phương pháp \textbf{One-vs-Rest} trong phân lớp đa lớp với $k$ lớp cần huấn luyện:
    \begin{enumerate}[label=\Alph*.]
        \item 1 bộ phân lớp.
        \item $k-1$ bộ phân lớp.
        \item $k$ bộ phân lớp.
        \item $k(k-1)/2$ bộ phân lớp.
    \end{enumerate}

    \item Trong Decision Tree, \textbf{Leaf node} chứa:
    \begin{enumerate}[label=\Alph*.]
        \item Điều kiện kiểm tra thuộc tính.
        \item Nhãn lớp (class label).
        \item Tập con của dữ liệu huấn luyện.
        \item Hệ số Gini Index.
    \end{enumerate}

    \item Information Gain $\text{Gain}(A) = \text{Info}(D) - \text{Info}_A(D)$ là:
    \begin{enumerate}[label=\Alph*.]
        \item Lượng thông tin cần thiết để phân lớp sau khi dùng thuộc tính A.
        \item Phần giảm entropy khi biết thuộc tính A (lớn hơn là tốt hơn).
        \item Độ đo Gini của phân hoạch.
        \item Xác suất thuộc tính A được chọn làm nút gốc.
    \end{enumerate}

    \item Thuật toán cây quyết định nào sử dụng \textbf{Gain Ratio}?
    \begin{enumerate}[label=\Alph*.]
        \item ID3.
        \item C4.5.
        \item CART.
        \item K-NN.
    \end{enumerate}

    \item Gini Index \textbf{nhỏ hơn} tương ứng với:
    \begin{enumerate}[label=\Alph*.]
        \item Partition có độ hỗn loạn cao hơn (impure).
        \item Partition thuần hơn (các đối tượng chủ yếu thuộc một lớp).
        \item Thuộc tính phân tách kém hơn.
        \item Entropy lớn hơn.
    \end{enumerate}

    \item CART (Classification and Regression Trees) đặc điểm là:
    \begin{enumerate}[label=\Alph*.]
        \item Cây đa phân (multiple branches per node).
        \item Cây nhị phân (binary split); sử dụng Gini Index.
        \item Cây nhị phân; sử dụng Information Gain.
        \item Cây đa phân; sử dụng Gain Ratio.
    \end{enumerate}

    \item Trong Naive Bayes, giả định \textbf{class conditional independence} nghĩa là:
    \begin{enumerate}[label=\Alph*.]
        \item Tất cả các lớp có xác suất bằng nhau.
        \item Các thuộc tính độc lập với nhau trong cùng một lớp.
        \item Các lớp độc lập với nhau.
        \item Thuộc tính liên tục và danh mục được xử lý giống nhau.
    \end{enumerate}

    \item \textbf{Laplace smoothing} trong Naive Bayes giải quyết:
    \begin{enumerate}[label=\Alph*.]
        \item Vấn đề dữ liệu thiếu (missing data).
        \item Vấn đề xác suất bằng 0 (zero-frequency problem).
        \item Vấn đề giả định độc lập bị vi phạm.
        \item Vấn đề chiều cao (high dimensionality).
    \end{enumerate}

    \item Công thức Laplace smoothing: $P(x_k \mid C_i) = \frac{\text{count} + 1}{|C_{i,D}| + m}$, trong đó $m$ là:
    \begin{enumerate}[label=\Alph*.]
        \item Số lớp trong tập dữ liệu.
        \item Số giá trị khác nhau của thuộc tính $A_k$.
        \item Số mẫu trong tập huấn luyện.
        \item Số thuộc tính của đối tượng.
    \end{enumerate}

    \item Trong ANN, ma trận trọng số $\Theta^{(j)}$ ánh xạ:
    \begin{enumerate}[label=\Alph*.]
        \item Từ lớp $j+1$ về lớp $j$.
        \item Từ lớp $j$ sang lớp $j+1$.
        \item Từ lớp đầu vào trực tiếp đến lớp đầu ra.
        \item Từ lớp ẩn sang lớp đầu vào.
    \end{enumerate}

    \item Kích thước của $\Theta^{(j)}$ trong ANN với $s_j$ nút tại lớp $j$ và $s_{j+1}$ nút tại lớp $j+1$ là:
    \begin{enumerate}[label=\Alph*.]
        \item $s_j \times s_{j+1}$
        \item $s_{j+1} \times (s_j + 1)$ (bao gồm bias)
        \item $(s_j + 1) \times s_{j+1}$
        \item $s_j \times (s_{j+1} + 1)$
    \end{enumerate}

    \item Trong Backpropagation, $\delta^{(l)}_j$ biểu diễn:
    \begin{enumerate}[label=\Alph*.]
        \item Giá trị kích hoạt của neuron $j$ tại lớp $l$.
        \item Sai số (error) tạo ra bởi neuron $j$ tại lớp $l$.
        \item Trọng số kết nối của neuron $j$ tại lớp $l$.
        \item Giá trị gradient của $J$ theo $\Theta^{(l)}_{ij}$.
    \end{enumerate}

    \item K-NN phân lớp đối tượng mới bằng cách:
    \begin{enumerate}[label=\Alph*.]
        \item Tính xác suất thuộc mỗi lớp theo Bayes.
        \item Bỏ phiếu đa số từ $k$ đối tượng gần nhất trong tập huấn luyện.
        \item Xây dựng cây quyết định từ tập huấn luyện.
        \item Tính sigmoid của $\theta^Tx$.
    \end{enumerate}

    \item Khi $k$ trong K-NN quá nhỏ, bộ phân lớp sẽ:
    \begin{enumerate}[label=\Alph*.]
        \item Ổn định hơn với nhiễu.
        \item Nhạy cảm hơn với nhiễu (overfitting).
        \item Có xu hướng phân lớp sai do chọn quá nhiều lớp.
        \item Không bị ảnh hưởng bởi $k$.
    \end{enumerate}

    \item Giá trị $k$ được gợi ý trong slide C4 là:
    \begin{enumerate}[label=\Alph*.]
        \item $k = 1$
        \item $k \leq \sqrt{|D|}$
        \item $k = 10$
        \item $k = |D| / 2$
    \end{enumerate}

    \item Precision được tính bằng:
    \begin{enumerate}[label=\Alph*.]
        \item $TP / (TP + FN)$
        \item $TP / (TP + FP)$
        \item $TN / (TN + FP)$
        \item $(TP + TN) / (TP + FP + FN + TN)$
    \end{enumerate}

    \item Recall được tính bằng:
    \begin{enumerate}[label=\Alph*.]
        \item $TP / (TP + FP)$
        \item $TP / (TP + FN)$
        \item $FP / (FP + TN)$
        \item $TN / (TN + FN)$
    \end{enumerate}

    \item F-score được tính bằng:
    \begin{enumerate}[label=\Alph*.]
        \item $(P + R) / 2$
        \item $2PR / (P + R)$
        \item $\sqrt{P \times R}$
        \item $P \times R$
    \end{enumerate}

    \item Với dataset 13 flows (9 BG và 4 FG), bộ phân lớp xác định 7 flows là BG (4 BG đúng và 3 FG sai). Precision và Recall của lớp BG lần lượt là:
    \begin{enumerate}[label=\Alph*.]
        \item P = 4/9, R = 4/7.
        \item P = 4/7, R = 4/9.
        \item P = 4/13, R = 4/7.
        \item P = 4/7, R = 4/13.
    \end{enumerate}

    \item Thuật toán cây quyết định nào sử dụng \textbf{Information Gain} và phù hợp với thuộc tính danh mục?
    \begin{enumerate}[label=\Alph*.]
        \item C4.5.
        \item CART.
        \item ID3.
        \item Naive Bayes.
    \end{enumerate}

    \item Trong ANN, lớp đầu vào $L=1$ chứa:
    \begin{enumerate}[label=\Alph*.]
        \item Kết quả dự đoán.
        \item Các đặc trưng đầu vào $x_1, \ldots, x_n$ (và bias $x_0 = 1$).
        \item Các trọng số $\Theta$.
        \item Giá trị delta (lỗi) của mỗi neuron.
    \end{enumerate}

    \item $\text{Info}(D) = -\sum_{i=1}^{m} p_i \log_2 p_i$ bằng \textbf{0} khi:
    \begin{enumerate}[label=\Alph*.]
        \item Tất cả đối tượng phân bổ đều giữa các lớp.
        \item Tất cả đối tượng thuộc cùng một lớp (partition thuần nhất).
        \item $m = 2$ lớp.
        \item $|D| = 1$.
    \end{enumerate}

    \item Bộ phân lớp nào \textbf{dễ giải thích nhất} (highest interpretability)?
    \begin{enumerate}[label=\Alph*.]
        \item ANN với nhiều hidden layers.
        \item K-NN.
        \item Decision Tree.
        \item Logistic Regression với nhiều biến.
    \end{enumerate}

    \item Ưu điểm chính của ANN so với Decision Tree là:
    \begin{enumerate}[label=\Alph*.]
        \item ANN dễ hiểu và giải thích hơn.
        \item ANN xử lý tốt hơn các bài toán phân lớp phi tuyến phức tạp với dữ liệu chiều cao.
        \item ANN cần ít dữ liệu huấn luyện hơn.
        \item ANN huấn luyện nhanh hơn với mọi loại dữ liệu.
    \end{enumerate}

    \item Tiêu chí \textbf{Scalability} trong đánh giá bộ phân lớp quan tâm đến:
    \begin{enumerate}[label=\Alph*.]
        \item Khả năng bộ phân lớp hiểu được kết quả.
        \item Khả năng xây dựng và cập nhật bộ phân lớp với tập dữ liệu rất lớn.
        \item Khả năng phân lớp đúng các đối tượng nhiễu.
        \item Tốc độ dự đoán trên tập test.
    \end{enumerate}

    \item Trong quy trình phân lớp 2 bước, bước \textbf{Evaluation} được thực hiện ở:
    \begin{enumerate}[label=\Alph*.]
        \item Chỉ ở bước Training.
        \item Ở đầu bước Classification (dùng testing dataset để đánh giá trước khi áp dụng).
        \item Sau khi đã phân lớp tất cả dữ liệu mới.
        \item Không cần thực hiện evaluation nếu training dataset đủ lớn.
    \end{enumerate}

    \item Phương trình $h_\theta(x) = g(\theta_0 + \theta_1 x_1^2 + \theta_2 x_2^2)$ tạo ra decision boundary dạng:
    \begin{enumerate}[label=\Alph*.]
        \item Đường thẳng.
        \item Đường tròn (hoặc elipse).
        \item Đường hyperbola.
        \item Không thể xác định.
    \end{enumerate}

    \item Trong K-fold cross-validation, tổng số lần một mẫu dữ liệu được dùng để test là:
    \begin{enumerate}[label=\Alph*.]
        \item $k$ lần.
        \item 1 lần.
        \item $k-1$ lần.
        \item $1/k$ lần.
    \end{enumerate}

    \item Trong Naive Bayes, $P(C_i) = |C_{i,D}| / |D|$ biểu diễn:
    \begin{enumerate}[label=\Alph*.]
        \item Likelihood (xác suất quan sát $X$ nếu lớp là $C_i$).
        \item Prior probability (xác suất tiên nghiệm của lớp $C_i$).
        \item Posterior probability (xác suất hậu nghiệm).
        \item Evidence (hằng số chuẩn hóa).
    \end{enumerate}

    \item ANN được lấy cảm hứng từ:
    \begin{enumerate}[label=\Alph*.]
        \item Hệ thống điều khiển robot.
        \item Cấu trúc và hoạt động của não người.
        \item Cấu trúc cây quyết định phân cấp.
        \item Lý thuyết tập mờ (Fuzzy sets).
    \end{enumerate}

    \item Khi Naive Bayes gặp vấn đề $P(x_k \mid C_i) = 0$ mà không dùng Laplace smoothing thì:
    \begin{enumerate}[label=\Alph*.]
        \item Chỉ ảnh hưởng đến thuộc tính đó.
        \item Toàn bộ $P(X \mid C_i) = 0$, loại bỏ hoàn toàn lớp $C_i$ khỏi xét.
        \item Tăng xác suất của các lớp khác.
        \item Không ảnh hưởng nếu có đủ thuộc tính khác.
    \end{enumerate}

    \item Trong đánh giá mô hình phân lớp, độ đo \textbf{Accuracy} có hạn chế khi:
    \begin{enumerate}[label=\Alph*.]
        \item Tập test quá nhỏ.
        \item Dữ liệu mất cân bằng lớp (class imbalance), ví dụ 99\% lớp A và 1\% lớp B.
        \item Số lớp lớn hơn 2.
        \item Bộ phân lớp sử dụng gradient descent.
    \end{enumerate}

\end{enumerate}

% ============================================================
\newpage
\section*{ĐÁP ÁN}
% ============================================================

\subsection*{Câu hỏi tự luận -- Hướng dẫn trả lời}

\begin{singlespace}
\begin{longtable}{|p{0.08\textwidth}|p{0.82\textwidth}|}
\hline
\textbf{Câu} & \textbf{Nội dung cần trình bày} \\
\hline
\endhead
1 & Classification: xây dựng hàm $f(X) = y$ từ dữ liệu có nhãn. Bước 1: Training (học mô hình từ $\langle X, y \rangle$). Bước 2: Classification (phân lớp dữ liệu mới). Supervised vì có nhãn. Khác clustering: clustering không có nhãn, tự tìm nhóm. \\
\hline
2 & Hồi quy tuyến tính: $h_\theta(X)$ có thể $> 1$ hoặc $< 0$ $\Rightarrow$ không phải xác suất. Sigmoid: $g(z) = 1/(1+e^{-z}) \in (0,1)$ $\Rightarrow$ đảm bảo $0 \leq h_\theta \leq 1$, có thể diễn giải là xác suất. \\
\hline
3 & $h_\theta(x) = g(\theta^Tx) \in (0,1)$. Ý nghĩa: $P(y=1|x;\theta)$. Quy tắc: dự đoán $y=1$ khi $h_\theta(x) \geq 0.5 \Leftrightarrow \theta^Tx \geq 0$; dự đoán $y=0$ khi ngược lại. \\
\hline
4 & Decision boundary là tập $\{\theta^Tx = 0\}$. Đường thẳng khi đặc trưng tuyến tính $\theta_0 + \theta_1 x_1 + \theta_2 x_2 = 0$. Đường cong khi thêm đặc trưng bậc cao: $\theta_0 + \theta_1 x_1^2 + \theta_2 x_2^2 = 0$ $\Rightarrow$ đường tròn. \\
\hline
5 & Dùng log-loss thay MSE vì: MSE với sigmoid cho $J$ không lồi (nhiều cực tiểu cục bộ). Log-loss: $y=1, h \to 0$: cost $\to \infty$ (phạt mạnh); $y=1, h=1$: cost $= 0$. Tương tự $y=0$. \\
\hline
6 & One-vs-Rest: với $k$ lớp, huấn luyện $k$ bộ phân lớp $h_\theta^{(i)}$, mỗi cái phân biệt lớp $i$ vs. tất cả còn lại. Dự đoán: lớp $i = \arg\max_i h_\theta^{(i)}(x)$. Hạn chế: không xét đồng thời tất cả lớp. \\
\hline
7 & Internal node: test thuộc tính (ví dụ: age $\leq 30$). Branch: kết quả test. Leaf: nhãn lớp. Thuật toán: greedy (chọn thuộc tính tốt nhất tại mỗi nút, không quay lại), divide-and-conquer, đệ quy, top-down. \\
\hline
8 & Information Gain: đo giảm entropy; có thể ưu tiên thuộc tính nhiều giá trị (bad). Gain Ratio: chuẩn hóa bằng SplitInfo $\Rightarrow$ khắc phục thiên vị của IG; dùng khi thuộc tính có nhiều giá trị. Gini: phân tách nhị phân, đơn giản hơn. \\
\hline
9 & Gain(age) = 0.940 - 0.694 = 0.246 bits. Nghĩa: Khi biết thuộc tính ``age'', lượng thông tin cần thêm để phân lớp giảm 0.246 bits. Để so sánh với Gain(income), Gain(student), Gain(credit\_rating) để chọn thuộc tính phân tách. \\
\hline
10 & Bayes: $P(H|X) = P(X|H)P(H)/P(X)$. Posterior: xác suất thuộc lớp sau khi quan sát $X$. Prior: xác suất lớp trước khi quan sát. Likelihood: xác suất quan sát $X$ nếu lớp là $H$. Phân lớp tối đa hóa posterior để chọn lớp có khả năng cao nhất. \\
\hline
11 & Giả định: $P(X|C_i) = \prod_k P(x_k|C_i)$ -- các thuộc tính độc lập trong mỗi lớp. Không thỏa khi: thuộc tính có tương quan (ví dụ: chiều cao và cân nặng). Hệ quả: xác suất tính toán không phản ánh đúng thực tế, có thể cho kết quả sai. \\
\hline
12 & Zero-frequency: $P(x_k|C_i) = 0$ khi giá trị $x_k$ chưa gặp trong lớp $C_i$ trong training $\Rightarrow$ $P(X|C_i) = 0$, loại bỏ lớp sai. Laplace: thêm 1 vào tử, $m$ vào mẫu. Thay thế: z-estimate: $P = (\text{count} + zP(x_k))/(|C_i| + z)$. \\
\hline
13 & Layer 1: input ($x_1, \ldots, x_n$). Hidden layers: xử lý phi tuyến. Layer L: output. $\Theta^{(j)}$: ma trận trọng số từ lớp $j$ sang $j+1$; kích thước $s_{j+1} \times (s_j + 1)$ (bias node thêm 1 vào $s_j$). \\
\hline
14 & Feedforward: $z^{(l+1)} = \Theta^{(l)} a^{(l)}$; $a^{(l+1)} = g(z^{(l+1)})$; thêm $a_0^{(l)}=1$ tại mỗi lớp. Bias node cần thiết để cho phép dịch chuyển (shift) decision boundary, tương tự intercept $\theta_0$ trong hồi quy. \\
\hline
15 & Backprop tính gradient $\partial J/\partial \Theta^{(l)}$. $\delta^{(L)} = a^{(L)} - y$ (lỗi output). $\delta^{(l)} = (\Theta^{(l)})^T\delta^{(l+1)} \cdot g'(z^{(l)})$ (lan truyền ngược). $\frac{\partial J}{\partial \Theta^{(l)}_{ij}} = a_j^{(l)}\delta_i^{(l+1)}$. Cập nhật $\Theta$ bằng gradient descent. \\
\hline
16 & XOR không phân tách tuyến tính được. ANN: có thể biểu diễn NAND $\to$ tổ hợp $\to$ XOR. ANN học các đặc trưng phi tuyến qua hidden layers, tạo ra biên phân lớp phức tạp. Logistic Regression chỉ tạo ra biên tuyến tính. \\
\hline
17 & K-NN: tính khoảng cách từ $X$ đến tất cả điểm trong $D$; chọn $k$ điểm gần nhất; dự đoán lớp có nhiều phiếu nhất. Chọn $k$: $k \leq \sqrt{|D|}$. Ưu: đơn giản. Nhược: tốn bộ nhớ, tính khoảng cách chậm khi $|D|$ lớn (lazy learner). \\
\hline
18 & Precision: trong số dự đoán là X, bao nhiêu phần đúng. Recall: trong số thực sự là X, bao nhiêu phần được phát hiện. Phát hiện gian lận: Recall quan trọng hơn (không bỏ sót gian lận thật). F-score cân bằng cả hai. \\
\hline
19 & Holdout: chia 1 lần, kết quả phụ thuộc cách chia. K-fold: chia $k$ lần, tính trung bình $\Rightarrow$ đáng tin cậy hơn vì ít phụ thuộc cách chia dữ liệu; tận dụng toàn bộ dữ liệu vừa train vừa test. \\
\hline
20 & Logistic Reg: accuracy trung bình, nhanh, ổn định, interpretable. Decision Tree: tốt với dữ liệu danh mục, nhanh, interpretable tốt. Naive Bayes: nhanh, ít dữ liệu, giả định độc lập. ANN: accuracy cao nhất (phi tuyến), chậm, kém interpretable. K-NN: accuracy tốt, chậm với data lớn, tốn bộ nhớ. \\
\hline
\end{longtable}
\end{singlespace}

\subsection*{Câu hỏi trắc nghiệm -- Đáp án}

\begin{singlespace}
\begin{longtable}{|c|c||c|c||c|c||c|c|}
\hline
\textbf{Câu} & \textbf{ĐA} & \textbf{Câu} & \textbf{ĐA} & \textbf{Câu} & \textbf{ĐA} & \textbf{Câu} & \textbf{ĐA} \\
\hline
\endhead
1 & B & 11 & B & 21 & B & 31 & C \\
\hline
2 & B & 12 & B & 22 & B & 32 & B \\
\hline
3 & A & 13 & B & 23 & B & 33 & B \\
\hline
4 & B & 14 & B & 24 & B & 34 & B \\
\hline
5 & B & 15 & B & 25 & C & 35 & B \\
\hline
6 & C & 16 & B & 26 & B & 36 & B \\
\hline
7 & B & 17 & B & 27 & B & 37 & B \\
\hline
8 & B & 18 & B & 28 & C & 38 & B \\
\hline
9 & B & 19 & B & 29 & C & 39 & B \\
\hline
10 & B & 20 & B & 30 & B & 40 & B \\
\hline
\end{longtable}
\end{singlespace}

\end{document}
